\documentclass[11pt,a4paper]{article}
% Preambulo comun de los reportes de Tecnologias Celulares.
% Cada main.tex hace: \documentclass ... % Preambulo comun de los reportes de Tecnologias Celulares.
% Cada main.tex hace: \documentclass ... % Preambulo comun de los reportes de Tecnologias Celulares.
% Cada main.tex hace: \documentclass ... \input{../comun/preambulo.tex}
\usepackage[utf8]{inputenc}
\usepackage[T1]{fontenc}
\usepackage[spanish,mexico]{babel}
\usepackage{graphicx}
\usepackage[margin=2.5cm]{geometry}
\usepackage{amsmath}
\usepackage{amssymb}
\usepackage{booktabs}
\usepackage{tabularx}
\usepackage{array}
\usepackage{float}
\usepackage{setspace}
\usepackage{parskip}
\usepackage{fancyhdr}
\usepackage{textcomp}
\usepackage{xcolor}
\usepackage{caption}
\usepackage{fvextra}
\usepackage[hidelinks]{hyperref}

\onehalfspacing
\setlength{\parskip}{6pt}
\sloppy
\setlength{\emergencystretch}{3em}
\setlength{\headheight}{40pt}
\captionsetup{font=small,labelfont=bf}
\setlength{\tabcolsep}{6pt}

\pagestyle{fancy}
\fancyhf{}
\fancyhead[L]{\includegraphics[width=1cm]{/home/yovick/repos/UAZ-IEI/Assets/UIE.png}}
\fancyhead[R]{\small Tecnologías Celulares}
\fancyfoot[C]{\thepage}
\renewcommand{\headrulewidth}{0.4pt}

\DefineVerbatimEnvironment{codigo}{Verbatim}{%
  breaklines=true,breakanywhere=true,fontsize=\small,frame=single,
  framesep=3pt,rulecolor=\color{gray!50}}

% Figura centrada con pie: \figura[ancho]{archivo}{leyenda}
\newcommand{\figura}[3][0.85]{%
  \begin{center}
    \includegraphics[width=#1\linewidth]{#2}\\[2pt]
    {\small #3}
  \end{center}}

% Portada: \portada{numero}{titulo de la practica}
\newcommand{\portada}[2]{%
\begin{titlepage}
  \thispagestyle{empty}
  \begin{center}
    \vspace*{1cm}
    \includegraphics[width=4.5cm]{/home/yovick/repos/UAZ-IEI/Assets/uaz.png}
    \vspace{0.8cm}

    {\large Universidad Autónoma de Zacatecas}\\[4pt]
    {\large Unidad Académica de Ingeniería Eléctrica}\\[4pt]
    {\normalsize Ingeniería Electrónica Industrial con especialidad en Telecomunicaciones}

    \vspace{1.8cm}
    {\LARGE\bfseries Práctica #1}\\[10pt]
    {\Large\bfseries #2}
  \end{center}

  \vspace{2.2cm}
  \begin{center}
    \begin{tabular}{rl}
      \textbf{Materia:} & Tecnologías Celulares\\
      \textbf{Docente:} & Dr. Víktor Iván Rodríguez Abdalá\\
      \textbf{Alumno:}  & Arnold Yovick Ríos Zaldívar\\
      \textbf{Fecha:}   & Noviembre de 2026\\
    \end{tabular}
  \end{center}
  \vfill
\end{titlepage}}

\usepackage[utf8]{inputenc}
\usepackage[T1]{fontenc}
\usepackage[spanish,mexico]{babel}
\usepackage{graphicx}
\usepackage[margin=2.5cm]{geometry}
\usepackage{amsmath}
\usepackage{amssymb}
\usepackage{booktabs}
\usepackage{tabularx}
\usepackage{array}
\usepackage{float}
\usepackage{setspace}
\usepackage{parskip}
\usepackage{fancyhdr}
\usepackage{textcomp}
\usepackage{xcolor}
\usepackage{caption}
\usepackage{fvextra}
\usepackage[hidelinks]{hyperref}

\onehalfspacing
\setlength{\parskip}{6pt}
\sloppy
\setlength{\emergencystretch}{3em}
\setlength{\headheight}{40pt}
\captionsetup{font=small,labelfont=bf}
\setlength{\tabcolsep}{6pt}

\pagestyle{fancy}
\fancyhf{}
\fancyhead[L]{\includegraphics[width=1cm]{/home/yovick/repos/UAZ-IEI/Assets/UIE.png}}
\fancyhead[R]{\small Tecnologías Celulares}
\fancyfoot[C]{\thepage}
\renewcommand{\headrulewidth}{0.4pt}

\DefineVerbatimEnvironment{codigo}{Verbatim}{%
  breaklines=true,breakanywhere=true,fontsize=\small,frame=single,
  framesep=3pt,rulecolor=\color{gray!50}}

% Figura centrada con pie: \figura[ancho]{archivo}{leyenda}
\newcommand{\figura}[3][0.85]{%
  \begin{center}
    \includegraphics[width=#1\linewidth]{#2}\\[2pt]
    {\small #3}
  \end{center}}

% Portada: \portada{numero}{titulo de la practica}
\newcommand{\portada}[2]{%
\begin{titlepage}
  \thispagestyle{empty}
  \begin{center}
    \vspace*{1cm}
    \includegraphics[width=4.5cm]{/home/yovick/repos/UAZ-IEI/Assets/uaz.png}
    \vspace{0.8cm}

    {\large Universidad Autónoma de Zacatecas}\\[4pt]
    {\large Unidad Académica de Ingeniería Eléctrica}\\[4pt]
    {\normalsize Ingeniería Electrónica Industrial con especialidad en Telecomunicaciones}

    \vspace{1.8cm}
    {\LARGE\bfseries Práctica #1}\\[10pt]
    {\Large\bfseries #2}
  \end{center}

  \vspace{2.2cm}
  \begin{center}
    \begin{tabular}{rl}
      \textbf{Materia:} & Tecnologías Celulares\\
      \textbf{Docente:} & Dr. Víktor Iván Rodríguez Abdalá\\
      \textbf{Alumno:}  & Arnold Yovick Ríos Zaldívar\\
      \textbf{Fecha:}   & Noviembre de 2026\\
    \end{tabular}
  \end{center}
  \vfill
\end{titlepage}}

\usepackage[utf8]{inputenc}
\usepackage[T1]{fontenc}
\usepackage[spanish,mexico]{babel}
\usepackage{graphicx}
\usepackage[margin=2.5cm]{geometry}
\usepackage{amsmath}
\usepackage{amssymb}
\usepackage{booktabs}
\usepackage{tabularx}
\usepackage{array}
\usepackage{float}
\usepackage{setspace}
\usepackage{parskip}
\usepackage{fancyhdr}
\usepackage{textcomp}
\usepackage{xcolor}
\usepackage{caption}
\usepackage{fvextra}
\usepackage[hidelinks]{hyperref}

\onehalfspacing
\setlength{\parskip}{6pt}
\sloppy
\setlength{\emergencystretch}{3em}
\setlength{\headheight}{40pt}
\captionsetup{font=small,labelfont=bf}
\setlength{\tabcolsep}{6pt}

\pagestyle{fancy}
\fancyhf{}
\fancyhead[L]{\includegraphics[width=1cm]{/home/yovick/repos/UAZ-IEI/Assets/UIE.png}}
\fancyhead[R]{\small Tecnologías Celulares}
\fancyfoot[C]{\thepage}
\renewcommand{\headrulewidth}{0.4pt}

\DefineVerbatimEnvironment{codigo}{Verbatim}{%
  breaklines=true,breakanywhere=true,fontsize=\small,frame=single,
  framesep=3pt,rulecolor=\color{gray!50}}

% Figura centrada con pie: \figura[ancho]{archivo}{leyenda}
\newcommand{\figura}[3][0.85]{%
  \begin{center}
    \includegraphics[width=#1\linewidth]{#2}\\[2pt]
    {\small #3}
  \end{center}}

% Portada: \portada{numero}{titulo de la practica}
\newcommand{\portada}[2]{%
\begin{titlepage}
  \thispagestyle{empty}
  \begin{center}
    \vspace*{1cm}
    \includegraphics[width=4.5cm]{/home/yovick/repos/UAZ-IEI/Assets/uaz.png}
    \vspace{0.8cm}

    {\large Universidad Autónoma de Zacatecas}\\[4pt]
    {\large Unidad Académica de Ingeniería Eléctrica}\\[4pt]
    {\normalsize Ingeniería Electrónica Industrial con especialidad en Telecomunicaciones}

    \vspace{1.8cm}
    {\LARGE\bfseries Práctica #1}\\[10pt]
    {\Large\bfseries #2}
  \end{center}

  \vspace{2.2cm}
  \begin{center}
    \begin{tabular}{rl}
      \textbf{Materia:} & Tecnologías Celulares\\
      \textbf{Docente:} & Dr. Víktor Iván Rodríguez Abdalá\\
      \textbf{Alumno:}  & Arnold Yovick Ríos Zaldívar\\
      \textbf{Fecha:}   & Noviembre de 2026\\
    \end{tabular}
  \end{center}
  \vfill
\end{titlepage}}


\begin{document}

\portada{2}{Generación y análisis de secuencias pseudoaleatorias (PN)}

\textbf{Script:} \texttt{practica2.py} (Python con NumPy; el manual pide MATLAB, pero el
algoritmo del LFSR y las correlaciones son idénticos)

\section{Objetivo}

Generar secuencias pseudoaleatorias de longitud máxima con registros de corrimiento con
retroalimentación lineal (LFSR) y analizar sus propiedades de balance, autocorrelación y
correlación cruzada.

\section{Fundamento teórico}

Un LFSR de $m$ etapas genera una secuencia binaria mediante la recurrencia lineal
\begin{equation}
  s_{n+m} = \bigoplus_{i \in T} s_{n+i} \pmod 2,
\end{equation}
donde $T$ es el conjunto de taps que define el polinomio característico. Si el polinomio
es primitivo, la secuencia recorre todos los estados no nulos del registro y su periodo es
\begin{equation}
  N = 2^m - 1.
\end{equation}
Para $m=5$ se usa el trinomio $x^5+x^2+1$, que da $N=31$; para $m=4,6,7$ se usan
$x^4+x+1$, $x^6+x+1$ y $x^7+x^3+1$, con longitudes 15, 63 y 127.

Una secuencia de longitud máxima cumple la propiedad de balance: en cada periodo hay
$2^{m-1}$ unos y $2^{m-1}-1$ ceros, es decir, la diferencia es exactamente uno. Su
autocorrelación periódica es de dos valores:
\begin{equation}
  R(\tau) =
  \begin{cases}
    1, & \tau = 0 \pmod N,\\
    -\dfrac{1}{N}, & \tau \neq 0 \pmod N.
  \end{cases}
\end{equation}
La correlación cruzada entre dos secuencias distintas debe mantenerse baja; para un par
preferido de secuencias-$m$ su valor máximo es del orden de $2^{(m+1)/2}/N$. Esa
propiedad es la que permite separar usuarios en CDMA: cada usuario multiplica sus datos
por un código casi ortogonal con los de los demás.

\section{Desarrollo experimental}

\subsection{Generación}

Se implementó el LFSR en forma de Fibonacci con estado inicial todos unos y salida por la
primera etapa. Para cada $m$ se generaron $N$ chips y se midió el periodo observado como
el primer desplazamiento en el que la secuencia se repite. La secuencia principal es la de
$m=5$ con $x^5+x^2+1$; para la correlación cruzada se generó una segunda secuencia de
$m=5$ con $x^5+x^4+1$.

\subsection{Análisis}

\begin{enumerate}
  \item Conversión bipolar $x_b[n]=2x[n]-1$, con valores $\{-1,+1\}$.
  \item Histograma para verificar el balance.
  \item Autocorrelación lineal normalizada (equivalente a \texttt{xcorr(x,'coeff')}) y
        autocorrelación periódica calculada como la transformada inversa de
        $|X(f)|^2$, que es la que exhibe la propiedad de dos valores.
  \item Correlación cruzada normalizada entre las dos secuencias de 31 chips.
\end{enumerate}

\section{Código}

\begin{codigo}
def lfsr(m, k, n=None, estado_inicial=None):
    n = n or 2**m - 1
    estado = list(estado_inicial or [1] * m)
    salida = np.empty(n, dtype=int)
    for i in range(n):
        salida[i] = estado[0]
        retro = estado[0] ^ estado[k]
        estado = estado[1:] + [retro]
    return salida

def autocorrelacion_periodica(x):
    n = len(x)
    r = np.fft.ifft(np.abs(np.fft.fft(x)) ** 2).real
    r = np.concatenate((r[-(n // 2):], r[: n // 2 + 1]))
    return r / r[n // 2]
\end{codigo}

\section{Resultados}

\subsection{Secuencia generada}

\figura{figs/secuencia.png}{Secuencia PN de 31 chips ($m=5$) y su versión bipolar.}

La secuencia contiene 16 unos y 15 ceros, exactamente el balance teórico
$2^{4}=16$ y $2^{4}-1=15$. El periodo observado es 31: la secuencia no se repite antes
de agotar el periodo.

\subsection{Distribución y balance}

\figura[0.62]{figs/histograma.png}{Histograma de la secuencia bipolar.}

\subsection{Autocorrelación}

\figura{figs/autocorrelacion.png}{Autocorrelación lineal (arriba) y periódica (abajo).}

La autocorrelación periódica presenta un pico unitario en $\tau=0$ y un valor constante
de $-1/31=-0.0323$ en todos los demás retardos. La autocorrelación lineal muestra el mismo
pico dominante, pero sus valores laterales oscilan porque la correlación se calcula sobre
una sola réplica finita y no sobre la secuencia periódica.

\subsection{Correlación cruzada}

\figura{figs/correlacion_cruzada.png}{Correlación cruzada entre las secuencias $x^5+x^2+1$ y $x^5+x^4+1$.}

El máximo de la correlación cruzada es \textbf{0.3226} y el valor medio absoluto
\textbf{0.0936}. Son valores bajos frente al pico de autocorrelación (1.0), lo que indica
que las dos secuencias son casi ortogonales.

\subsection{Longitudes para distintos $m$}

\begin{tabular}{@{}cccccc@{}}
  \toprule
  $m$ & Polinomio & $N$ teórica & $N$ observada & Unos & Ceros \\
  \midrule
  4 & $x^4+x+1$   & 15  & 15  & 8  & 7  \\
  5 & $x^5+x^2+1$ & 31  & 31  & 16 & 15 \\
  6 & $x^6+x+1$   & 63  & 63  & 32 & 31 \\
  7 & $x^7+x^3+1$ & 127 & 127 & 64 & 63 \\
  \bottomrule
\end{tabular}

\section{Análisis de resultados}

\begin{itemize}
  \item \textbf{¿Por qué las secuencias PN son útiles en espectro disperso?} Porque son
        determinísticas (el receptor puede reproducirlas) pero se comportan como ruido:
        tienen espectro ancho y plano, y su autocorrelación impulsiva permite recuperar la
        sincronía y los datos mediante correlación.
  \item \textbf{¿Qué características presenta la autocorrelación?} Un pico agudo en el
        origen y un nivel residual constante de $-1/N$. Cuanto mayor es $N$, más pequeño
        es el residuo y más se parece la autocorrelación a un impulso.
  \item \textbf{¿Por qué la correlación cruzada es importante en CDMA?} Porque los códigos
        de distintos usuarios se suman en el mismo canal; si la correlación cruzada es
        baja, el receptor de un usuario rechaza la energía de los demás y solo ve la
        interferencia residual.
  \item \textbf{¿Cómo cambia la longitud de la secuencia al aumentar $m$?} Crece
        exponencialmente: $N=2^m-1$, de 15 a 127 al pasar de $m=4$ a $m=7$. El número de
        estados posibles se duplica con cada etapa y el residuo de autocorrelación se
        reduce a la mitad.
\end{itemize}

\section{Cuestionario}

\begin{enumerate}
  \item \textbf{¿Qué es una secuencia PN?} Una secuencia binaria periódica generada por un
        proceso determinístico que imita las propiedades estadísticas del ruido: balance
        de unos y ceros, autocorrelación impulsiva y baja correlación cruzada.
  \item \textbf{¿Qué es un LFSR?} Un registro de corrimiento de $m$ etapas cuya entrada es
        una combinación lineal (XOR) de algunas de sus salidas. Con un polinomio primitivo
        recorre todos los estados no nulos y produce una secuencia de longitud máxima.
  \item \textbf{¿Qué significa secuencia de longitud máxima?} Que su periodo es el máximo
        posible para un registro de $m$ etapas, $N=2^m-1$, es decir, que pasa por todos
        los estados no nulos antes de repetirse.
  \item \textbf{¿Por qué las secuencias PN son determinísticas pero parecen aleatorias?}
        Porque las genera una recurrencia fija y reproducible, pero su periodo largo, su
        balance y su espectro plano hacen que un observador que no conoce el polinomio ni
        el estado inicial no pueda predecirlas ni distinguirlas de ruido.
  \item \textbf{¿Cuál es la ventaja de una autocorrelación impulsiva?} Permite sincronizar
        el receptor y separar la señal útil del ruido y de otras señales: al correlacionar
        con la réplica local, solo el retardo correcto produce un pico grande.
  \item \textbf{¿Por qué se utilizan secuencias PN en sistemas celulares?} Para dispersar
        el espectro (DSSS/CDMA), identificar celdas y usuarios, sincronizar tramas y
        aleatorizar la transmisión. En LTE, por ejemplo, las secuencias Zadoff-Chu y los
        códigos Gold cumplen funciones análogas de sincronización e identificación.
\end{enumerate}

\section{Conclusiones}

Se implementó un LFSR de Fibonacci y se generaron secuencias de longitud máxima para
$m=4,5,6,7$, con periodos 15, 31, 63 y 127 y balance $2^{m-1}$ unos contra
$2^{m-1}-1$ ceros. La secuencia de 31 chips presentó autocorrelación periódica de dos
valores (1 en el origen y $-0.0323$ fuera de él) y una correlación cruzada máxima de
0.3226 con otra secuencia del mismo periodo. Estas propiedades son la base de los
sistemas de espectro disperso y del acceso múltiple por código.

\end{document}
